%% chapters/behavior_modeling/ch4_abms.tex
%% 第4章　アクティビティモデルと離散連続モデル

\section{アクティビティモデル}\label{sec:activity_models}
本節では，行動モデルが記述する選択行動を単独のトリップから一日の活動スケジュール全体へと拡張する．
前節までは交通手段選択や経路選択といった個別のトリップに関する離散選択を扱ってきたが，
そもそもトリップは活動の派生需要であり，個人の行動をより包括的に記述するには活動スケジュール全体を対象とする必要がある．
このような視点に立つモデルを総称して\textbf{アクティビティモデル}（Activity-Based Model; ABM）と呼ぶ．

本節ではまずアクティビティモデルの基本的な考え方を導入し，
続いて活動種類の選択と時間配分を同時に扱う\textbf{離散連続モデル}について解説する．

\subsection{活動ベースアプローチ}

従来の交通需要予測で広く用いられてきた四段階推定法は，生成・分布・分担・配分の各段階を順に解くことで交通需要を推定する．
しかしこの枠組みではトリップを独立した単位として扱うため，一日の中で活動がどのように連鎖し，トリップがいかにして生じるかという行動の文脈を十分に捉えることができない．

\textbf{活動ベースアプローチ}（Activity-Based Approach）は，トリップそのものではなく，個人の活動スケジュールを分析の基本単位とする枠組みである．
その基本的な考え方は以下の通りである．

\begin{description}
  \item[派生需要としてのトリップ]
    個人は通勤，買い物，余暇などの活動に参加するためにトリップを行う．
    すなわちトリップは活動の派生需要であり，活動スケジュールが決まればトリップのパタンも決まる．
  \item[一日の活動スケジュール]
    個人の行動は「活動種類・開始時刻・終了時刻・活動場所」の系列として記述される．
    この系列を\textbf{活動経路}（Activity Path）と呼ぶ．
  \item[制約の考慮]
    活動スケジュールは時間的制約（一日は24時間），空間的制約（移動には時間がかかる），
    および個人属性による制約（自動車免許の有無など）のもとで決定される．
\end{description}

活動ベースアプローチの利点は，活動間の連鎖や時間配分の制約を明示的に扱える点にある．
例えば，通勤帰りに買い物に立ち寄る行動（トリップチェイン）や，勤務時間の変更が余暇活動に与える影響などは，トリップ単位の分析では捉えにくいが，活動スケジュール全体を対象とすることで自然に記述できる．

\subsection{アクティビティモデルの構造}

アクティビティモデルは，活動スケジュールの生成過程を記述するモデルであり，その構造は大きく以下の二つに分類される．

\subsubsection{階層的選択モデル}
Bowman and Ben-Akiva \cite{BOWMAN20011} に代表される階層的選択モデルは，活動スケジュールの決定を階層的な離散選択の連鎖として記述する．
具体的には，まず一日の活動パタン（在宅・通勤・買い物など）を選択し，
次に主活動の時間帯を選択し，さらに目的地や交通手段を選択するという階層構造を仮定する．
各階層はNested Logitなどの離散選択モデルによって記述され，
上位の選択が下位の選択に影響を与える構造を持つ．

この枠組みは理論的整合性と解釈可能性に優れるが，
階層構造の設定は分析者の事前知識に依存するため，
複雑な活動パタンの表現には限界がある．

\subsubsection{逐次的選択モデル}
Zimmermann et al. \cite{ZIMMERMANN2018273} は，Recursive Logitモデル \cite{Fosgerau2013,Oyama2017dRLeng} を活動スケジューリングに適用し，
活動の逐次的な選択過程を記述するモデルを提案した．
このアプローチでは，各時点において次に行う活動を逐次的に選択していくことで活動スケジュールを生成する．
前節で紹介した経路非列挙型モデルと同様の発想であり，
活動経路の列挙を必要とせず，選択肢集合の爆発を回避できるという利点がある．

\subsection{日間変動とマルチデイモデル}

個人の活動スケジュールは日ごとに変動する．
例えば，毎日の通勤パタンは概ね規則的であるが，買い物や余暇活動は曜日や直近の活動履歴によって変化する．
このような日間変動を捉えるために，\textbf{マルチデイモデル}が提案されている．

代表的なアプローチとして，活動ニーズ（Activity Needs）に基づくモデルがある \cite{ARENTZE2009251, ARENTZE2011447}．
このモデルでは，各活動種類に対して「前回の実施からの経過時間」に応じて活動ニーズが蓄積し，
ニーズが閾値を超えたときに活動が発生すると仮定する．
また，Habib \cite{habib2008modelling} は活動ニーズの充足を効用として定式化し，効用最大化の枠組みの中で活動発生を記述するモデルを提案した．

これらのモデルは日間変動の構造を明示的に表現できる一方で，
閾値やニーズ蓄積関数の設定が分析者の事前知識に依存すること，
また活動場所を同時に扱う場合に次元の呪いが顕著になることが課題として挙げられる．

\subsection{離散連続モデル}\label{ssec:discrete_continuous}

前項までのモデルは主に「どの活動を行うか」という離散選択に焦点を当てていた．
しかし，活動スケジューリングにおいては活動種類の選択（離散選択）と同時に，
各活動への時間配分（連続選択）も決定される．
このような離散選択と連続選択を同時に扱うモデルを\textbf{離散連続モデル}（Discrete-Continuous Model）と呼ぶ．

\subsubsection{Corner Solution の問題}

離散連続モデルの必要性を理解するために，まず\textbf{Corner Solution}の問題を考える．

個人が一日の時間を複数の活動種類に配分する状況を考えよう．
全部で $K$ 種類の活動が存在するとき，個人は各活動 $k$ に対して時間 $t_k \ge 0$ を配分し，
$\sum_{k=1}^{K} t_k = T$（$T$ は総利用可能時間）という制約のもとで効用を最大化する．

ここで重要な点は，多くの活動種類に対して配分時間が $t_k = 0$ となることである．
例えば，ある日に買い物をしなかった個人にとって，買い物への時間配分は $t_k = 0$ である．
このように最適解が定義域の境界上に位置する状況をCorner Solutionと呼ぶ．

通常の連続最適化では内点解を前提とするため，Corner Solutionを適切に扱うことができない．
離散連続モデルはこの問題に対処するために開発されたものであり，
「どの活動を行うか（離散選択）」と「行う活動にどれだけ時間を配分するか（連続選択）」を
統一的な枠組みで記述する．

\subsubsection{MDCEV モデル}

離散連続モデルの中で最も広く用いられているのが，Bhat \cite{BHAT2008274} によって提案された\textbf{Multiple Discrete-Continuous Extreme Value}（MDCEV）モデルである．

MDCEVモデルでは，個人の効用関数を以下のように定義する．
\begin{equation}\label{eq:mdcev_utility}
  U(\bm{t}) = \sum_{k=1}^{K} \frac{\gamma_k}{\alpha_k} \psi_k \left[ \left( \frac{t_k}{\gamma_k} + 1 \right)^{\alpha_k} - 1 \right]
\end{equation}
ここで $\bm{t} = (t_1, \dots, t_K)$ は各活動への時間配分ベクトル，$\psi_k = \exp(\bm{\beta}^\top \bm{x}_k + \varepsilon_k)$ は活動 $k$ の基本効用（ベースライン効用），$\gamma_k > 0$ は活動 $k$ の消費単位を調整するパラメータ，$\alpha_k \le 1$ は飽和パラメータである．

この効用関数の特徴は以下の通りである．
\begin{description}
  \item[飽和効果]
    $\alpha_k < 1$ のとき，活動 $k$ への時間配分が増加するにつれて限界効用が逓減する．
    これにより，一つの活動に全ての時間を配分するのではなく，複数の活動に時間を分散させる行動が自然に表現される．
  \item[Corner Solutionの許容]
    $t_k = 0$ が最適解となることを許容する構造を持つ．
    すなわち，基本効用 $\psi_k$ が十分に小さい活動には時間が配分されないという行動を記述できる．
  \item[閉じた形の選択確率]
    誤差項 $\varepsilon_k$ にi.i.d. Gumbel分布を仮定することで，
    どの活動の組合せが選択されるかの確率が閉じた形で得られる．
\end{description}

MDCEVモデルは活動種類と時間配分の同時決定を統一的に扱えるため，
アクティビティモデルにおける時間配分の記述に広く応用されている \cite{BHAT2008274, habib2011random}．

\subsubsection{Tobit モデル}

離散連続モデルのもう一つの古典的な枠組みとして\textbf{Tobit モデル}がある．
Tobit モデルは，潜在変数 $t_k^*$ が閾値を超えた場合にのみ正の値が観測されるという構造を持つ．

\begin{equation}\label{eq:tobit}
  t_k = \max(0, \, t_k^*), \quad t_k^* = \bm{\beta}^\top \bm{x}_k + \varepsilon_k
\end{equation}

Tobit モデルは構造が単純で推定が容易である一方，
複数の活動種類間の相互作用や時間制約を明示的に扱うことが難しい．
このため，活動スケジューリングのように複数の離散連続選択が同時に生じる問題には，
MDCEVモデルの方がより適した枠組みであるといえる．
